\section{Reconfiguration-Phase Accountability Races}
\label{sec:reconf}

A reviewer's first question about any single-shard, single-epoch attack is: what happens when
validators rotate among shards? Sharded systems reassign committees and hand off shard state and
cross-shard artifacts at epoch boundaries~\cite{zamani2018rapidchain,kokoris2018omniledger}. This
section extends the model to that boundary. Its thesis is deliberately narrow:
\emph{reconfiguration does not eliminate the accountability race; it relocates it.} Whether
relocation helps the attacker or the defender depends entirely on how old-epoch certified artifacts
are handed off, and we give the exact condition (Propositions~\ref{prop:reconf-safe}
and~\ref{prop:reconf-residual}).

\paragraph{Three events that must be distinguished.}
We do \emph{not} claim that validator rotation lets bribed validators escape slashing. Robust
PoS/BFT systems separate three events that a careless analysis conflates:
(1) \emph{shard rotation}---a validator leaves shard $A$ and joins shard $B$;
(2) \emph{validator exit}---a validator stops validating; and
(3) \emph{stake withdrawal}---a validator recovers bonded stake and is no longer slashable.
A validator that merely rotates out of a shard should remain slashable for old-epoch equivocation,
because the evidence---validator identity, shard ID, epoch, height/round, and the two conflicting
signed values---is self-contained and can be enforced by a global staking/beacon layer after
rotation (Counterexample~\ref{cex:persistent}). The attack's reconfiguration angle is therefore
\emph{not} ``rotation defeats slashing''; it is whether a \emph{value-bearing certified artifact}
created by the old committee can execute in the next epoch before old-epoch evidence is enforced.

\paragraph{Variables.}
Let $\tau$ be the epoch length and $a\in[0,\tau]$ the attack time within the epoch, with
$r=\tau-a$ the remaining time before nominal reconfiguration. Let $U$ be the unbonding/withdrawal
delay after validator exit, $\Hreconf(a)$ the extra accountability delay caused by committee handoff,
state transfer, beacon processing, or proof routing across epochs, $\Tstate$ the time for the new
committee to obtain and verify shard state, $\Treceipt$ the time for old-epoch receipts/locks/
withdrawal certificates/outbox messages to become executable in the next epoch, and $\Tevid$ the time
for old-epoch equivocation evidence to become enforceable against globally bonded stake. Finally,
$q(a)=\Pr[\text{evidence enforced before stake unslashable}]$ generalizes the enforcement
probability $q$ of \S\ref{sec:cost} to attack time $a$.

\paragraph{Extended profit model.}
The single-epoch model of \S\ref{sec:design} becomes, as a function of attack time,
\begin{equation}
\mathbb{E}[\Pi(a)] \;=\; \psucc(a)\,\pwin(a)\,\Eval \;-\; B(a) \;-\; \Cop,
\label{eq:profit-a}
\end{equation}
with $B(a)=(F{+}1)\,q(a)\,(s{+}R)$. The window term depends on the handoff design. For
\emph{receipt-carryover} designs, where the old-epoch artifact survives into epoch $e{+}1$ and
executes there,
\begin{equation}
\pwin(a)=\Pr\big[\,\Tsettle^{\mathrm{ce}} < \Tacc + \Hreconf(a)\,\big],
\label{eq:pwin-carry}
\end{equation}
where the cross-epoch settlement time $\Tsettle^{\mathrm{ce}}=\Treceipt+\Tstate+\Twd$ (with $\Twd$
the destination withdrawal) must beat $\Tevid$: the handoff penalty $\Hreconf(a)$ \emph{widens} the
effective accountability deadline, so the artifact need only beat $\Tevid$. For
\emph{hard-cutoff} designs, where the old committee stops at the boundary and no value-bearing
artifact may execute afterward,
\begin{equation}
\pwin(a)=\Pr\big[\,\Tsettle < \min(\Tacc,\,r)\,\big].
\label{eq:pwin-cut}
\end{equation}
As $a\to\tau$ we have $r\to0$, so last-block attacks become \emph{useless} under a hard cutoff unless
the malicious artifact was already certified and settles within the shrinking $r$. The two regimes
pull $\mathbb{E}[\Pi(a)]$ in opposite directions near the boundary; Figure~\ref{fig:reconf-profit}
shows the resulting curves.

\subsection{What Survives Reconfiguration}
\label{sec:survives}

Ordinary mempool transactions are \emph{not} the main risk. A pending transaction can be dropped,
retried, re-gossiped, or carried into the next epoch without ever being consensus-final; nothing
value-bearing has been certified. The dangerous objects are old-epoch \emph{certified} artifacts:
finalized shard headers, state roots, receipt/outbox roots, pending incoming receipts, cross-shard
locks, withdrawal certificates, and bridge-relevant certificates. The attack matters precisely when
an old committee creates a certified value-bearing artifact that remains \emph{executable} after
reconfiguration but before old-epoch equivocation evidence is enforced. We therefore restate the
race in handoff terms:

\begin{quote}
\emph{Reconfiguration does not eliminate the accountability race; it relocates it. Even if the old
committee stops producing blocks, certified receipts, locks, and state roots may survive into the
next epoch. The relevant residual window is not the time left to produce ordinary blocks, but the
time during which old-epoch artifacts remain executable before old-epoch equivocation evidence is
enforced.}
\end{quote}

\subsection{A Taxonomy of Reconfiguration Designs}
\label{sec:taxonomy}

Table~\ref{tab:taxonomy} classifies the handoff designs by what they do at the boundary and the
condition under which each is safe or dangerous. Figure~\ref{fig:reconf-fsm} renders the handoff as a
state machine with the safe (gated) and unsafe (ungated) paths.

\begin{table*}[t]
\centering
\caption{Reconfiguration designs and their effect on the attack. ``Safe'' is the condition under
which $\pwin\to0$ for old-epoch value-bearing artifacts (Proposition~\ref{prop:reconf-safe});
``dangerous'' is the residual-vulnerability condition (Proposition~\ref{prop:reconf-residual}).}
\label{tab:taxonomy}
\small
\begin{tabularx}{\textwidth}{@{}l X X X X@{}}
\toprule
Design & At the epoch boundary & Effect on attack & Safe condition & Dangerous condition \\
\midrule
A. Hard cutoff + mempool carryover & Old committee stops; unincluded txns retried/re-gossiped & Last-block attack weak unless the malicious tx was already certified & Value-bearing artifacts cannot execute after the cutoff & A receipt was certified before the cutoff and settles within $r$ \\
B. Drain phase & Protocol stops accepting new cross-shard txns and finishes pending locks/receipts before rotation & Strong if drained receipts execute before accountability & Pending receipts held through the accountability window & Drained receipts execute immediately at boundary \\
C. Receipt/outbox handoff & Old-epoch receipts/outbox imported in epoch $e{+}1$ & The most important case: artifact survives and executes next epoch & Import is accountability-gated (held until evidence window elapses) & Imported receipts executed immediately on import \\
D. Old/new committee overlap & Old committee signs final checkpoint; new committee signs import checkpoint & Depends on whether the new committee checks for conflicts & New committee rejects/quarantines conflicting old certificates & New committee blindly imports one old root \\
E. Global beacon/control-layer checkpoint & Beacon records shard headers, receipt roots, assignments, slashing evidence & Preserves accountability across rotation & Beacon quarantines old receipt roots until evidence can land & Beacon finalizes old receipt roots before evidence is enforced \\
F. Rolling/asynchronous reconfiguration & Validators rotate gradually; no single last-block cliff & No cutoff, but attacker optimizes $a$ via $q(a),\Tstate,\Tevid$ & Old-epoch artifacts stay gated throughout the roll & Some artifacts execute during the roll before evidence \\
\bottomrule
\end{tabularx}
\end{table*}

\begin{figure}[t]
\centering
\includegraphics[width=0.96\linewidth]{fig5_reconf_fsm.pdf}
\caption{Epoch-boundary handoff as a state machine. The safe (accountability-gated) path quarantines
old-epoch artifacts until evidence can land ($\pwin=0$); the unsafe (ungated) path executes them
immediately, reopening the race.}
\label{fig:reconf-fsm}
\end{figure}

\subsection{State Handoff Under Committee Rotation}
\label{sec:handoff}

New validators must obtain shard state before they can validate. Mechanisms include a full snapshot
transfer verified against a state root; a checkpoint plus block replay; stateless validation with
witnesses/Merkle proofs; old/new committee checkpoint overlap; a global beacon/control-layer shard
header; and erasure-coded or data-availability-backed state reconstruction. Crucially, the new
committee needs more than account balances. It needs the (1) latest finalized shard header,
(2) state root, (3) receipt/outbox root, (4) pending incoming receipts, (5) pending cross-shard
locks, (6) withdrawal certificates, (7) validator-assignment metadata, (8) old-epoch
equivocation/slashing evidence, and optionally (9) mempool transactions. Of these, mempool state is
\emph{not} consensus-critical---it can be reconstructed or dropped---whereas receipt roots, locks,
finalized headers, and equivocation evidence \emph{are}. A handoff that imports value-bearing roots
(3)--(6) without first importing and checking evidence (8), and without gating execution behind the
source accountability window, is exactly the dangerous case of Table~\ref{tab:taxonomy}.

\subsection{Counterexamples: When the Attack Does Not Apply}
\label{sec:counterexamples}

Precise threat models name their boundaries. The following cases make the reconfiguration angle
reviewer-resistant by stating when it fails.

\begin{counterexample}[Persistent slashability after rotation]\label{cex:persistent}
A validator rotates from shard $A$ to shard $B$, but its stake remains globally bonded. The
equivocation evidence (identity, shard ID, epoch, height/round, conflicting signed values) is
self-contained, so a global slashing contract or beacon layer can slash after rotation. Shard
rotation therefore does \emph{not} imply accountability escape; $q(a)$ stays high and $B(a)$ does not
fall.
\end{counterexample}

\begin{counterexample}[Accountability-gated receipt carryover]\label{cex:gated}
Old-epoch receipts carry into epoch $e{+}1$ but remain \emph{pending} through a challenge/
accountability window. If equivocation evidence appears, the receipt root is quarantined and the
validators slashed. This enforces $\Tsettle\ge\Tacc$ and gives $\pwin=0$ (Figure~\ref{fig:reconf-profit},
flat line): the attack is killed regardless of $a$.
\end{counterexample}

\begin{counterexample}[Cross-shard freeze during reconfiguration]\label{cex:freeze}
The protocol allows local transactions but freezes cross-shard receipt execution during state
handoff. Last-block cross-shard attacks fail because no value-bearing artifact can settle across the
boundary.
\end{counterexample}

\begin{counterexample}[Rolling reconfiguration]\label{cex:rolling}
There is no single point where the old committee disappears, so there is no last-block escape to
exploit. The attack must be modeled continuously; with old-epoch artifacts gated throughout the roll,
$\pwin$ never opens (Figure~\ref{fig:reconf-profit}, rolling curve stays bounded).
\end{counterexample}

These cases do not weaken the paper; they sharpen its threat model.

\subsection{Two Propositions}
\label{sec:reconf-props}

\begin{proposition}[Reconfiguration-safe handoff]\label{prop:reconf-safe}
Suppose a sharded protocol satisfies all of:
\emph{(i)} old-epoch equivocation evidence remains valid and enforceable after committee rotation;
\emph{(ii)} validator stake remains bonded and slashable for at least the maximum evidence delay
$\Tevid$;
\emph{(iii)} old-epoch receipts, locks, and withdrawal certificates are not executable until the
corresponding source-shard accountability window has elapsed; and
\emph{(iv)} the new committee verifies imported state, receipt roots, and pending artifacts against a
canonical checkpoint and rejects/quarantines conflicting old-epoch certificates.
Then committee reconfiguration does not increase the adversary's settlement-before-accountability
probability for old-epoch artifacts; in particular, accountability-gated old-epoch artifacts have
$\pwin=0$.
\end{proposition}
\begin{proof}
See Appendix~\ref{app:proofs}.
\end{proof}

\begin{proposition}[Reconfiguration-residual vulnerability]\label{prop:reconf-residual}
Suppose instead that the protocol allows old-epoch certified receipts, locks, or withdrawal artifacts
to become executable in epoch $e{+}1$ before old-epoch equivocation evidence can be enforced
(violating (iii)), and that reconfiguration either increases accountability latency by $\Hreconf(a)$
with $\Hreconf$ nondecreasing in $a$, or decreases $q(a)$ near the epoch boundary. Then there is a
boundary region $[\tau-\delta,\tau]$ on which $\mathbb{E}[\Pi(a)]$ is increasing in $a$; attacking
near the end of the epoch strictly increases expected adversarial profit relative to a mid-epoch
attack.
\end{proposition}
\begin{proof}
See Appendix~\ref{app:proofs}.
\end{proof}

\begin{figure}[t]
\centering
\includegraphics[width=\linewidth]{fig6_reconf_profit.pdf}
\caption{Expected profit $\mathbb{E}[\Pi(a)]$ across the epoch for each handoff design
(\S\ref{sec:eval}, $N=31$, $V=60s$). Hard cutoff (A) collapses as $a\to\tau$ (residual time $r\to0$);
unsafe carryover (C) and reconfiguration-delay designs \emph{rise} toward the boundary
(Proposition~\ref{prop:reconf-residual}); accountability-gated carryover stays flat and negative
($\pwin=0$, Proposition~\ref{prop:reconf-safe}); rolling reconfiguration (F) has no cliff.}
\label{fig:reconf-profit}
\end{figure}
